\chapter{Nachweisumfang und wissenschaftliche Aussagen}
\label{vv:scope}\index{Verifikation}\index{Validierung}
Dieser Band beschreibt den nachprüfbaren Stand der linearen FEM Toolbox bis M8
und des generischen M9-Unsicherheitsworkflows. Er begründet keine experimentelle
Validierungsfreigabe. Die am 10. Oktober 2026 genehmigte Neuplanung von M9
trennt die Software- und Dokumentationsfreigabe v1.4.0 vom weiterhin offenen
Messvergleich in Issues 63 und 65. Das ursprüngliche Messdatengate ist damit
kein v1.4.0-Releasegate mehr; seine wissenschaftlichen Anforderungen gelten
unverändert für einen späteren experimentellen Fall.

\section{Drei unterschiedliche Fragen}
\textbf{Code-Verifikation} fragt, ob die Implementierung die spezifizierten
Gleichungen korrekt löst. Elementmatrizen, Handrechnungen, Starrkörper- und
Affinpatches, Bilanzgleichungen und hergestellte Lösungen liefern dazu
unterschiedliche, einander ergänzende Referenzen. Ein zweiter Solver mit
denselben assemblierten Matrizen ist eine Backendprüfung, aber keine unabhängige
Prüfung des Stoffgesetzes.

\textbf{Lösungsverifikation} untersucht den numerischen Fehler einer konkreten
Rechnung: Netz- und Zeitschrittfolgen, algebraische Genauigkeit und geeignete
Zielgrößen. Ein kleines Residuum kann bei einem zu groben Netz mit einem großen
Diskretisierungsfehler zusammenfallen. Ein bestandener analytischer Stabtest
überträgt seine Fehlergrenze nicht auf beliebige Bauteile.

\textbf{Experimentelle Validierung} vergleicht ein physikalisches Modell mit
unabhängigen Messungen für einen festgelegten Verwendungszweck. Dafür fehlen
gegenwärtig angenommene Messdaten mit vollständigem Proben-, Material-,
Randbedingungs- und Unsicherheitsvertrag. Synthetische Beispiele erfüllen diese
Aufgabe nicht. Die unten dokumentierten negativen Suchbefunde sind ein
Beschaffungsbefund, kein Beweis weltweiter Nichtverfügbarkeit.

\section{Unabhängigkeit der Referenzen}
Die stärksten elementaren Referenzen entstehen außerhalb der Produktionsroutine:
analytische Tensorrotation, Simplexmomente, direkt integrierte Polynome und
bekannte eindimensionale Lösungen. Bei der Method of Manufactured Solutions
wird ein glattes Feld gewählt und die zugehörige rechte Seite durch Einsetzen
in die Differentialgleichung abgeleitet \cite{sandia-mms}. Die Referenz soll
nicht durch erneuten Aufruf desselben Operators erzeugt werden.

Für quadratische Tetraeder nutzt die Konvergenzauswertung unter anderem eine
unabhängige Duffy-Integration mit $6\times6\times6=216$ Punkten. Ein Test gegen
einen gespeicherten Referenzwert erkennt Regressionen, kann aber einen Fehler
in der ursprünglichen Referenz konservieren. Deshalb ersetzen solche Tests
die analytischen und integralen Nachweise nicht.

\section{Leseregel für die Nachweismatrix}
Die folgenden 17 Gruppen verknüpfen jede unterstützte Physik bzw. wichtige
Querschnittsfunktion mit ihrem Vertrag und tatsächlich vorhandenen Tests.
Die Matrix wird aus \path{Doc/latex/verification-map.json} erzeugt; der
Build prüft Vollständigkeit und Existenz aller Referenzen. Sie bezeichnet den
\emph{Prüfumfang}, nicht ein unabhängig zertifiziertes Produkt.
Ein lokaler Lauf ohne PETSc überspringt native Tests und kann diese daher
nicht als bestanden melden. Native Ergebnisse gehören zum zugehörigen
CI-Commit und seiner realen oder komplexen PETSc-Umgebung.

% Generated by scripts/verification_artifacts.py; review sources before regeneration.
\subsection{Stationäre Wärme}
\label{sec:vv-heat}
Tet4/Tet10; SciPy; native PETSc-Parität separat. Affine Felder, Fluss-/Leistungsbilanz, vier Netze; P1 L2/H1 2/1, P2 3/2.
Fachvertrag: \path{Doc/45_TET10_HEAT_INTEGRATION.md}.
\begin{itemize}
\item \path{tests/patch/test_heat_patch.py}
\item \path{tests/convergence/test_heat_convergence.py}
\item \path{tests/convergence/test_tet10_heat_operator.py}
\end{itemize}
\subsection{Transiente Wärme}
\label{sec:vv-transient_heat}
Tet4/Tet10; Backward Euler. Kapazität/Robin/Ereignisse; räumliche Normen und Zeitordnung 1 getrennt.
Fachvertrag: \path{Doc/33_TRANSIENT_HEAT.md}.
\begin{itemize}
\item \path{tests/convergence/test_transient_heat_convergence.py}
\item \path{tests/convergence/test_tet10_transient_heat.py}
\item \path{tests/unit/test_transient_heat.py}
\end{itemize}
\subsection{Elektrostatik}
\label{sec:vv-electrostatics}
Tet4/Tet10; SciPy; native PETSc-Parität separat. Elektrisches Feld und Leiterladung, Energie, Gauge und vier Netze.
Fachvertrag: \path{Doc/46_TET10_ELECTROSTATICS.md}.
\begin{itemize}
\item \path{tests/patch/test_electrostatics_patch.py}
\item \path{tests/convergence/test_electrostatics_convergence.py}
\item \path{tests/convergence/test_tet10_electrostatics.py}
\end{itemize}
\subsection{Elastizität}
\label{sec:vv-elasticity}
Tet4/Tet10; linear/stationär. Starrkörper und Affinpatch, Zug und Biegung, integrierte Energie, Kraft und Moment, vier Netze.
Fachvertrag: \path{Doc/47_TET10_ELASTICITY.md}.
\begin{itemize}
\item \path{tests/patch/test_elasticity_patch.py}
\item \path{tests/convergence/test_elasticity_convergence.py}
\item \path{tests/convergence/test_tet10_elasticity.py}
\end{itemize}
\subsection{Thermomechanik}
\label{sec:vv-thermomechanics}
Tet4/Tet10; einseitiges T nach u. Freie/behinderte Eigenverzerrung, getrennte Bilanzen und vier Netze.
Fachvertrag: \path{Doc/48_TET10_THERMOMECHANICS.md}.
\begin{itemize}
\item \path{tests/patch/test_thermomechanics_patch.py}
\item \path{tests/convergence/test_thermomechanics_convergence.py}
\item \path{tests/convergence/test_tet10_thermomechanics.py}
\end{itemize}
\subsection{Piezoelektrizität}
\label{sec:vv-piezoelectric}
Tet4/Tet10; stationär. d33 direkt/invers, Kurzschluss/Leerlauf, Reziprozität, Enthalpie und vier Netze.
Fachvertrag: \path{Doc/49_TET10_PIEZOELECTRIC.md}.
\begin{itemize}
\item \path{tests/patch/test_piezo_element_patch.py}
\item \path{tests/patch/test_piezo_solver_patch.py}
\item \path{tests/patch/test_electrodes.py}
\item \path{tests/regression/test_piezo_d33_benchmark.py}
\item \path{tests/convergence/test_tet10_piezoelectric.py}
\end{itemize}
\subsection{Thermopiezoelektrizität}
\label{sec:vv-thermopiezoelectric}
Tet4/Tet10; T nach (u,phi). Materialgrenze, Elektroden, thermischer Schalter, Feldnormen und getrennte Bilanzen.
Fachvertrag: \path{Doc/50_TET10_THERMOPIEZOELECTRIC.md}.
\begin{itemize}
\item \path{tests/convergence/test_thermopiezo_model_convergence.py}
\item \path{tests/convergence/test_tet10_thermopiezoelectric.py}
\item \path{tests/unit/test_thermopiezoelectric.py}
\end{itemize}
\subsection{Piezo-Moden}
\label{sec:vv-modal}
Tet4/Tet10; begrenzte dichte Eigenanalyse. Analytischer Stab, Massennormierung, rekonstruierte Schur-Lösung, vier Netze.
Fachvertrag: \path{Doc/51_TET10_MODES.md}.
\begin{itemize}
\item \path{tests/unit/test_piezo_modal.py}
\item \path{tests/convergence/test_piezo_modal_convergence.py}
\item \path{tests/convergence/test_tet10_modal.py}
\end{itemize}
\subsection{Piezo-Newmark}
\label{sec:vv-newmark}
Tet4/Tet10; beta=1/4, gamma=1/2. Ein-DOF-Oszillator, Energie, Zeitordnung 2; bewegte Ränder/Ereignisse.
Fachvertrag: \path{Doc/52_TET10_NEWMARK.md}.
\begin{itemize}
\item \path{tests/unit/test_piezo_newmark.py}
\item \path{tests/unit/test_tet10_newmark.py}
\end{itemize}
\subsection{Piezo-Harmonik}
\label{sec:vv-harmonic}
Tet4/Tet10; komplexe Spitzenphasoren. Analytischer Stab, DC, Nullamplitude/Phase, Resonanzablehnung, vier Netze.
Fachvertrag: \path{Doc/53_TET10_HARMONIC.md}.
\begin{itemize}
\item \path{tests/unit/test_piezo_harmonic.py}
\item \path{tests/convergence/test_tet10_harmonic_convergence.py}
\end{itemize}
\subsection{Geometrie und Materialien}
\label{sec:vv-geometry}
Affine Tet4; geradkantige Tet10/Triangle6. Simplexmomente, Polynomreproduktion, Permutationen, Tensor-/Voigtrotation, Negativfälle.
Fachvertrag: \path{Doc/41_TET10_GEOMETRY.md}.
\begin{itemize}
\item \path{tests/unit/test_tet10_geometry.py}
\item \path{tests/unit/test_tet10_operators.py}
\item \path{tests/unit/test_piezo_material.py}
\item \path{tests/unit/test_material_regions.py}
\end{itemize}
\subsection{Assembly und Backends}
\label{sec:vv-backends}
SciPy direkt; PETSc real/complex direkt/FieldSplit. Originalreaktionen, feldweise Residuen und native Parität; keine MPI-/Leistungszusage.
Fachvertrag: \path{Doc/39_PETSC_FIELDSPLIT.md}.
\begin{itemize}
\item \path{tests/unit/test_assembly.py}
\item \path{tests/unit/test_linear_backend.py}
\item \path{tests/petsc/test_physics.py}
\item \path{tests/petsc/test_fieldsplit.py}
\item \path{tests/petsc/test_convergence.py}
\end{itemize}
\subsection{Wärmeindikator}
\label{sec:vv-indicator}
Nur stationäre Tet4. Residuen, Flächenanteile, Dörfler-Markierung; keine garantierte Fehlerschranke.
Fachvertrag: \path{Doc/43_HEAT_INDICATOR.md}.
\begin{itemize}
\item \path{tests/unit/test_heat_indicator.py}
\item \path{tests/convergence/test_heat_indicator_convergence.py}
\end{itemize}
\subsection{Studien und Kennwerte}
\label{sec:vv-studies}
Stationäre Model-Schema-1-Fälle. Provenienz, Wiederaufnahme, Fehlerstatus, Tabellenhash und Zieleinheiten.
Fachvertrag: \path{Doc/60_STUDY_REPORTS.md}.
\begin{itemize}
\item \path{tests/unit/test_studies.py}
\item \path{tests/unit/test_study_reports.py}
\item \path{tests/regression/test_study_recipes.py}
\end{itemize}
\subsection{Polynomiale Surrogate}
\label{sec:vv-surrogates}
Reelle skalare Ziele; Grad 1/2. Analytische Interaktionen, unabhängige Holdouts, Ridge, Rang/Leckage/Extrapolation.
Fachvertrag: \path{Doc/61_SURROGATE_MODELS.md}.
\begin{itemize}
\item \path{tests/unit/test_surrogates.py}
\item \path{tests/regression/test_surrogate_example.py}
\end{itemize}
\subsection{Unsicherheitsworkflow}
\label{sec:vv-uncertainty}
Reelle skalare Ziele; stationäre Studien. Analytische Momente, Korrelation/Seeds, lokale Ableitungen, Fehler-/Wiederaufnahmefälle.
Fachvertrag: \path{Doc/67_UNCERTAINTY_WORKFLOW.md}.
\begin{itemize}
\item \path{tests/unit/test_uncertainty.py}
\item \path{tests/regression/test_uncertainty_study.py}
\end{itemize}
\subsection{Genauigkeit und Kosten}
\label{sec:vv-cost}
Tet4/Tet10; SciPy/PETSc direkt. 72 historische Läufe, unabhängige Normen; unveränderte Budgets und unerreichte Ziele.
Fachvertrag: \path{Doc/54_ACCURACY_COST_COMPARISON.md}.
\begin{itemize}
\item \path{tests/regression/test_accuracy_cost.py}
\item \path{tests/petsc/test_native_accuracy_cost.py}
\end{itemize}


\section{Gültigkeitsgrenzen}
Die Diskretisierung verwendet affine Tet4 und geradkantige Tet10; gekrümmte
isoparametrische Zellen sind kein belegter Umfang. Die Stoffgesetze sind linear.
Thermomechanik und Thermopiezoelektrizität übertragen zunächst Temperatur in
die mechanischen bzw. elektromechanischen Gleichungen. Die hier dokumentierte
Kopplung enthält keine vollständig rückgekoppelte nichtlineare Thermodynamik.
Der Wärmeindikator gilt nur für stationäre Tet4 und liefert keine allgemein
garantierte Fehlerschranke.

Moden und Newmark sind durch die dokumentierten Größenlimits begrenzt.
Harmonik bezeichnet eine einzelne undämpfte Frequenz mit komplexen
Spitzenphasoren; Resonanzablehnung ist ein negatives Akzeptanzkriterium.
Native PETSc-Parität und FieldSplit-Tests belegen kleine unterstützte Fälle.
Sie begründen weder MPI-Skalierbarkeit noch einen allgemeinen Geschwindigkeitsvorteil.
